\section{The light-delayed belief model}
\label{sec:model}

\subsection{The knowledge gate}
The change to the model is one gate on knowledge. A settled star is treated as an
omnidirectional beacon that begins emitting in the year it is settled. A probe deciding
its next target at star $f$ in year $Y$ believes a distant star $i$ is settled only once
that beacon's light has had time to arrive:
\begin{equation}
\text{settled}(i) + \frac{\mathrm{dist}(f, i)}{c} \le Y . \label{eq:gate}
\end{equation}
Under the perfect-information assumption of the source model, (\ref{eq:gate}) collapses to
``settled at any non-negative year,'' and the model reproduces our perfect-information
baseline bit-for-bit; the light-speed term is what introduces stale views. When two probes
both believe an unsettled star is the best target they both fly to it; the loser arrives to
find it taken, its trip wasted, and must re-target from the even more stale picture it holds
on arrival. This is availability chosen over consistency: rather than wait for a single
agreed map that finite light-speed makes impossible, each probe acts now on its local view
\cite{lamport-1978}. Stars carry velocity magnitudes that feed the slingshot boost but are
held fixed in position, following the source model, so $\mathrm{dist}(f,i)$ in
(\ref{eq:gate}) is static and the beacon geometry does not shift under the probe.

\subsection{The three flight policies}
We carry three named policies, following the slingshot studies we extend
\cite{forgan-papadog-kitching-2013,nicholson-forgan-2013}. \emph{Powered
nearest-neighbour} flight moves under the probe's own power at a cruise speed of
$3\times10^{-5}\,c$ to the nearest unsettled star. \emph{Nearest-star slingshot} flight
takes the same nearest target but gains speed from gravitational slingshots off the stars
it passes. \emph{Maximum-boost slingshot} flight instead selects, among the nearest
believed-unsettled stars, the target that yields the largest velocity gain, deliberately
reaching past near neighbours to distant, fast stars. The three span the axis that matters
here: powered flight is slow and local, nearest-star slingshot is fast and local, and
maximum-boost slingshot is fast and far. The effective speeds the policies reach in the
simulation, which set the lag ratio below and let the reader check the reported penalties,
are given with the results in Table~\ref{tab:mech}: a powered cruise near $9$ km/s against
slingshot means near $2{,}700$ km/s (nearest) and $3{,}900$ km/s (maximum-boost).

\subsection{The lag ratio, and what it can and cannot say}
How much the gate can matter on a single hop is governed by a dimensionless ratio: the
information delay across the hop divided by the travel time across the same hop. For a hop
of length $d$ the delay is $d/c$ and the travel time is $d/v$, so the ratio is
\begin{equation}
\Lambda \;=\; \frac{d/c}{d/v} \;=\; \frac{v}{c}, \label{eq:lambda}
\end{equation}
the probe's speed in units of $c$, and it is the \emph{same on every hop}: the length $d$
cancels. A slow probe ($\Lambda \ll 1$) outpaces its own ignorance, so stale views barely
arise; at the powered cruise speed $\Lambda \approx 3\times10^{-5}$ is negligible. A fast
probe carries a large $\Lambda$ and routinely decides before the relevant news can reach it.
So $\Lambda$ gates \emph{whether} the delay can bite at all, and it does so correctly: the
effect appears only once probes are fast, which is the slingshot regime. But precisely
because $\Lambda$ is hop-length-independent it cannot say \emph{where} the cost lands, and
Section~\ref{sec:results} shows the penalty is decided by hop length, not by $\Lambda$. We
therefore keep $\Lambda$ for its gating role and demote it as a predictor of the penalty's
size.

\subsection{Experimental design and parameters}
The field is a box of stars at a uniform density of one star per cubic parsec, the
convention of the source model. Each settled star launches two offspring (the replication
branching factor); this is what makes the field fill exponentially rather than as a single
roving chain, and it sets the swarm density that governs how often two probes contend for
one star. Three properties keep the measurement honest. First, knowledge is evaluated only
at the decision star at the moment of decision; news a probe's path happens to cross in
mid-flight is ignored. This deliberately undercounts what a probe could know, so the measured
waste is a conservative upper bound on the true penalty. Second, a probe that loses a race
re-targets, but only up to a fixed number of times before it is retired as wasted; this cap
bounds pathological bounce chains late in the fill, is bookkeeping rather than physics, and
the results are insensitive to it. Third, the whole simulation is a pure, seeded, fixed-step
fold: fixing the seed fixes the star field, every arrival, and every target choice, so a run
with the light-speed gate and a run without it can be compared on the \emph{identical} galaxy.
That paired design is what lets us attribute a difference in fill-time to the delay alone
rather than to the luck of the star layout.

The uniform density is easy to over-read, so one clarification. Changing the \emph{value} of
the uniform density rescales every inter-star distance by a common factor, which multiplies
both the travel time and the light-travel time on every hop equally; $\Lambda$ and, in the
continuum limit, the relative penalty are left unchanged (a pure geometric rescaling of the
absolute clock). We confirm this holds in practice: re-running at densities down to the
sparser solar-neighbourhood value near $0.14$ stars per cubic parsec \cite{recons-census}
keeps the median penalty within a few points of its value at unit density (about
\mbox{27-29\%} for nearest-star, about \mbox{46-52\%} for maximum-boost), the residual scatter
being the fixed timestep, not the density. A \emph{non-uniform}, clumpier distribution is a different matter, and we
treat it as a limitation (Section~\ref{sec:limitations}): by lengthening the long-range hops
that carry the penalty it is expected to move the relative penalties, not merely the absolute
timescale.

\subsection{Baseline validation}
The paired comparison is only as good as its perfect-information baseline, so we check that
baseline at two levels. By construction the ``instant'' regime is the $c\to\infty$ limit of
(\ref{eq:gate}), and it reproduces the default perfect-information fold bit-for-bit. Against
Nicholson and Forgan's published model \cite{nicholson-forgan-2013} we reproduce their two
qualitative findings: gravitational slingshots explore far faster than powered flight, and
their surprising result that aiming for the nearest star beats aiming for the largest boost
on total exploration time. In our fold, on a $400$-star field, nearest-star slingshot fills
in about $75{,}000$ years against about $305{,}000$ for maximum-boost, while maximum-boost
reaches the higher peak speed. We do not reproduce their roughly hundred-fold slingshot
speed-up quantitatively: at the fixed timestep the boosted hops are shorter than one step and
are quantized, so the measured speed-up is about twenty-fold. Lowering the timestep recovers
more of it; the qualitative ordering, on which this paper's conclusions rest, is
timestep-robust.
